\subsection{A large local Klein character and an exact cancellation at one}

\begin{proposition}[P-F40: the Klein normal representation does not supply an unsigned global Laurent bound]
\label{prop:kleinphase}
Let $n\ge4$ with $3\nmid n$, and put
\[
 Q=2^n-(-1)^n,\qquad P=Q/3,\qquad
 K_n=\sum_{i=0}^{n-1}x_i^2x_{i+1},\qquad x_n=x_0.
 \tag{P-F40a}
\]
The cubic $K_n$ is smooth. Set $w_i=(-2)^i$ and
$\rho=e^{2\pi i/P}$. The matrix
\[
 \beta=\operatorname{diag}(\rho^{w_0},\ldots,\rho^{w_{n-1}})
 \tag{P-F40b}
\]
fixes $K_n$ exactly, belongs to $SL_n$, and has order $P$.
Let $W$ be the space of cubic forms with this substitution
convention and let
\[
 S=W/(\mathfrak{sl}_n\cdot K_n),\qquad D=\binom n3.
 \tag{P-F40c}
\]
This is the normal representation to the $SL_n$ orbit, restricted
to the actual cyclic stabilizer subgroup $\Gamma=\langle\beta\rangle$.
It satisfies
\[
 \dim S=D+1,\qquad\dim S^\beta=1,\qquad
 \det(\beta\mid S)=1.
 \tag{P-F40d}
\]
For each $a$ coprime to $P$, define
\[
 C_a=\frac1{\det(1-\beta^a\mid S^*_{\rm mov})}.
\]
Then
\[
                 C_{-a}=\overline{C_a}=(-1)^D C_a.
 \tag{P-F40e}
\]
For the standard root in (P-F40b), the absolute value nevertheless
has the same large asymptotic scale as the chain calculation:
\[
                         \log_2|C_1|=\frac{n^4}{24}+O(n^3).
 \tag{P-F40f}
\]
The implied constant is absolute as $n$ ranges over $3\nmid n$.

For the cyclic normal-representation model, with ordinary
polynomial degree on $\operatorname{Sym}(S^*)$, Molien averaging is
\[
 \operatorname{Hilb}_{\operatorname{Sym}(S^*)^\Gamma}(t)
       =\frac1P\sum_{a=0}^{P-1}
                       \frac1{\det(1-t\beta^a\mid S^*)}.
 \tag{P-F40g}
\]
Every summand with $(a,P)=1$ has at most a simple pole at one,
with normalized coefficient $C_a/P$.
If $n\equiv3\pmod4$, the pair of summands indexed by $a$ and
$-a$ is holomorphic at one: its possible simple poles cancel
exactly. Thus the entire primitive-order part of (P-F40g)
is holomorphic at one in this parity. If $n\not\equiv3\pmod4$,
the two normalized coefficients add to $2C_a/P$, which is real;
its sign is specified below and is not assumed positive.
Already for $n=8$ the signs are mixed: $C_1<0<C_3$.

These statements concern an actual stabilizer representation
and its cyclic finite-quotient model. No identification of
(P-F40g), or any of its summands, with the global pole-at-one
polynomial $B_n$ of P-F34 is asserted. In particular (P-F40f)
is neither a lower nor an upper bound on the actual $U_n$.
\end{proposition}

\begin{proof}
The partial derivatives of $K_n$ are
\[
                 2x_i x_{i+1}+x_{i-1}^2
\]
with cyclic indices. If one coordinate vanishes at a common
zero of these derivatives, the preceding coordinate also
vanishes, and induction around the cycle gives the origin.
If no coordinate vanishes, multiplying all $n$ equations gives
$2^n(\prod_i x_i)^2=(-1)^n(\prod_i x_i)^2$, which is impossible
in characteristic zero. This proves smoothness.

For comparison with the known family, $K_n$ is the type-$K_n$
polynomial of Zheng, Section~2.2. Lemma~3.1, printed p.~5,
identifies its diagonal projective automorphism group as
cyclic of order $Q/3$. We prove the particular lift and
representation properties needed here directly.
Let $r=-2$ and $S_0=\sum_{i=0}^{n-1}r^i=(1-r^n)/3$.
Then $S_0=\pm P$. The relations
$2w_i+w_{i+1}=0$ for $i<n-1$ and
$2w_{n-1}+w_0=1-r^n=3S_0$ show that $\beta$ fixes
all $n$ terms of $K_n$. Also
$\det\beta=\rho^{S_0}=1$, and its first eigenvalue is $\rho$,
so its order is $P$.
One may obtain this lift by cubing the exact-form diagonal
matrix with a primitive $Q$th root; that original matrix has
determinant of order three. This distinction between linear
and projective orders is essential here.

The hypothesis $3\nmid n$ implies $3\nmid P$. Indeed
$(-2)^n=(1-3)^n\equiv1-3n\pmod9$, so
$(1-(-2)^n)/3\equiv n\pmod3$.
Every $w_i$ is congruent to one modulo three.
Consequently, if $P$ divides a difference $w_i-w_j$, then
$Q=3P$ divides it. But for distinct indices its absolute
value is at most $2^{n-1}+2^{n-2}<Q$.
The eigenvalues of $\beta$ are therefore distinct.

We next determine all fixed cubic monomials, rather than
assuming that the displayed support is the whole fixed space.
A cubic monomial with indices $i,j,k$ is fixed exactly when
$P\mid b$, where $b=r^i+r^j+r^k$. Since $3\mid b$ and
$(3,P)=1$, this is equivalent to $Q\mid b$.
The bound $|b|\le3\cdot2^{n-1}<2Q$ leaves only
$b=0,Q,-Q$.
If $b=0$, the least index occurring cannot have multiplicity
one or three, as is seen by dividing by its power of two.
It has multiplicity two, and
$2r^i+r^j=0$ forces $j=i+1$.
These are the $n-1$ nonwrapping terms of $K_n$.

If $b=\pm Q$, it is odd. Thus the index zero occurs an odd
number of times. It cannot occur three times, since
$Q\ge15$, so it occurs once and the other two indices
$j,k$ are positive. For the target $b=1-r^n$, one has
$r^j+r^k=-r^n$. If $j\ne k$, the two-adic valuation of
the left side is $\min(j,k)<n$; if $j=k$, the equation
forces $j=k=n-1$. This is the wrapping monomial
$x_{n-1}^2x_0$. For the other target $b=r^n-1$,
one would have $r^j+r^k=r^n-2$, whose two-adic valuation
is one. Equal indices would give valuation at least two.
For unequal indices the smaller must be one, and subtracting
$r^1=-2$ would force the larger to be $n$, which is not
an available index. There are no other fixed monomials.
Thus $\dim W^\beta=n$.

We justify the orbit dimension without assuming the full
finite stabilizer group has been classified. Let
$R=\mathbb C[x_0,\ldots,x_{n-1}]$ and
$\mathfrak m=(x_0,\ldots,x_{n-1})$.
The $n$ quadratic partial derivatives have common zero only
at the origin. In the regular local ring $R_{\mathfrak m}$
they generate an ideal of colength finite, so they are a
regular sequence. The exact Koszul complex says that any
relation among them has each coefficient in their generated
ideal. In particular a relation with linear coefficients
has every coefficient in $\mathfrak m^2R_{\mathfrak m}$,
and hence every such linear coefficient is zero.
A Lie algebra element fixing $K_n$ gives precisely such a
relation. Thus the map
$\mathfrak{sl}_n\to W$, $X\mapsto X\cdot K_n$,
is injective. The regular-sequence and Koszul facts used
here are Stacks Project Lemmas~10.106.3, 10.104.2 and
15.31.2, Tags~00NQ, 02JN and 062F, respectively.

The quotient sequence in (P-F40c) is $\Gamma$-equivariant.
Finite-group invariants over $\mathbb C$ are exact by
averaging. As $\beta$ has distinct eigenvalues,
$\dim(\mathfrak{sl}_n)^\beta=n-1$. Therefore
\[
 \dim S=\binom{n+2}{3}-(n^2-1)=D+1,
 \qquad \dim S^\beta=n-(n-1)=1.
\]
The determinant of the cubic representation at any matrix
$B$ equals $(\det B)^{(n+1)(n+2)/2}$: by symmetry each
coordinate eigenvalue occurs with total exponent
$3\binom{n+2}{3}/n$ in the product of all cubic monomial
eigenvalues. In particular it is one at $\beta$.
The adjoint determinant is also one, since off-diagonal
weights occur in inverse pairs and diagonal weights are one.
Taking determinants in the exact sequence gives
$\det(\beta\mid S)=1$. Its fixed part has determinant one,
so both the moving part and its dual have determinant one.
All these conclusions hold for $\beta^a$ when $(a,P)=1$;
passing to such a power does not change any fixed eigenspace.

Let $\alpha_1,\ldots,\alpha_D$ be the nontrivial eigenvalues
of $\beta^a$ on $S^*$. Their product is one, and
\[
 \prod_{j=1}^D(1-\alpha_j^{-1})
     =(-1)^D\prod_{j=1}^D(1-\alpha_j).
\]
They are roots of unity, so inverse also means complex
conjugate. Taking reciprocals proves (P-F40e).
For an exact sign description, write
$\alpha_j=e^{2\pi i\theta_j}$ with $0<\theta_j<1$ and put
$h_a=\sum_j\theta_j\in\mathbb Z$; integrality follows
from their product being one. Then
\[
 C_a=(-1)^{h_a}i^D |C_a|.
 \tag{P-F40h}
\]
Indeed $1-e^{2\pi i\theta}=-2i\,e^{\pi i\theta}
\sin(\pi\theta)$ and the sine is positive in this interval.
The integer $h_a$ is computable directly from the cubic
weights:
\[
 h_a=\sum_{0\le i\le j\le k<n}
       \left\{-\frac{a(w_i+w_j+w_k)}P\right\}
                   -\frac{n(n-1)}2,
 \tag{P-F40i}
\]
where braces denote fractional part. Fixed cubic weights
have fractional part zero; each off-diagonal adjoint inverse
pair contributes one to its age. Additivity in the exact
sequence proves (P-F40i).
For even $D$ this gives the exact real sign
$(-1)^{D/2+h_a}$. It supplies no positivity assertion for
sums over different $a$.

The large absolute value has an independent elementary
proof. The moving determinant quotient gives
\[
 |C_1|=
 \frac{\prod_{i\ne j}|1-\rho^{w_i-w_j}|}
 {\prod_{i\le j\le k,\ P\nmid w_i+w_j+w_k}
                    |1-\rho^{w_i+w_j+w_k}|}.
 \tag{P-F40j}
\]
Dualizing does not change these absolute values.
For largest index $k\le n-6$, the moving cubic counts
are exactly $N_0=1$ and $N_k=(k^2+3k)/2$ for $k\ge1$,
as in P-F34. Every corresponding nonzero integer weight
satisfies $2^{k-2}\le|b|\le3\cdot2^k$.
Since $P\ge2^{n-2}$, these weights satisfy $|b|<P/2$.
The sine inequalities give
\[
       \frac{2^k}{P}\le|1-\rho^b|
                   \le\frac{6\pi\,2^k}{P}.
\]
Because $\log_2P=n+O(1)$, the logarithm of each bulk
factor is $k-n+O(1)$ with an absolute constant.
Their total main contribution to $\log_2|C_1|$ is
\[
 \sum_{k=0}^{n-6}(n-k)N_k=\frac{n^4}{24}+O(n^3),
\]
with the exact sum already displayed in P-F34.
The only remaining cubic factors have their largest
index among the last five indices. There are $O(n^2)$
of these. All nontrivial $P$th-root factors lie between
$4/P$ and two, so each of these logarithms is $O(n)$.
The $n(n-1)$ adjoint factors satisfy the same estimate.
The total error is therefore $O(n^3)$, proving (P-F40f).
Dividing by $P$ would change this logarithm by only $O(n)$.

Finally, averaging the trace projector over $\Gamma$ on
every symmetric degree gives (P-F40g). For a diagonal
operator, the generating trace on its symmetric algebra
is the product of the geometric series
$(1-t\alpha)^{-1}$, which proves the formula directly.
When $(a,P)=1$, its one-dimensional fixed subspace gives
\[
 \frac1{\det(1-t\beta^a\mid S^*)}
       =\frac{\phi_a(t)}{1-t},\qquad
 \phi_a\text{ holomorphic at }1,\quad\phi_a(1)=C_a.
\]
If $D$ is odd, (P-F40e) makes
$\phi_a(1)+\phi_{-a}(1)=0$; the numerator of the paired
term is divisible by $1-t$ as a germ. The pair is
therefore holomorphic at one. If $D$ is even, its
normalized coefficient is $2C_a$ instead.
The integer $D=n(n-1)(n-2)/6$ is odd exactly when
$n\equiv3\pmod4$: among three consecutive factors,
its numerator has two-adic valuation one exactly
in that case, while the denominator has valuation one.
This proves every cancellation assertion.

For completeness, mixed signs in even codimension can be checked
exactly in the small case $n=8$, rather than left as a hypothetical
possibility. Here $P=85$ and $D=56$. For a residue $b$, write
$\mathcal O_b=\{b(-2)^j\bmod85:0\le j<8\}$, with residues in
$\{0,\ldots,84\}$. The moving dual cubic and adjoint representations
decompose into the following orbits. The last two columns give the
sum of the orbit's residues divided by $85$, before and after
multiplication by three.
\[
\begin{array}{c|rrrr}
 b&\text{cubic multiplicity}&\text{adjoint multiplicity}
       &\sum_{u\in\mathcal O_b}u/85
       &\sum_{u\in\mathcal O_b}(3u\bmod85)/85\\\hline
 1&1&0&4&4\\
 3&3&2&4&4\\
 6&3&2&4&4\\
 7&1&0&3&4\\
 9&2&1&4&3\\
15&2&1&4&4\\
18&2&1&4&5
\end{array}
\]
Here is a direct way to verify the whole finite table. Rotation
of the eight variable indices acts freely on cubic monomials,
because a nontrivial rotational stabilizer has even order,
which cannot divide the total monomial degree three.
The fifteen monomial orbits have representatives: the cube
$(0,0,0)$; the seven triples $(0,0,d)$ for $1\le d\le7$;
and the seven squarefree triples
\[
 (0,1,2),(0,1,3),(0,1,4),(0,1,5),(0,1,6),(0,2,4),(0,2,5).
\]
The last list follows by writing the three positive cyclic gaps
with sum eight, up to cyclic rotation. Reducing the negative
weight sum of each representative modulo $85$ gives exactly
the displayed cubic multiplicities; $(0,0,1)$ is the fixed
orbit and is omitted. The seven off-diagonal adjoint orbits
have representatives $(-2)^d-1$, $1\le d\le7$, giving the
second multiplicity column. Each displayed residue orbit has
eight distinct elements. Its residues, or their multiples by
three, give the two elementary integer sums in the table.
Subtracting the adjoint multiplicities leaves each of the seven
displayed orbits once in $S^*_{\rm mov}$. Thus
\[
 h_1=4+4+4+3+4+4+4=27,\qquad
 h_3=4+4+4+4+3+4+5=28.
\]
Formula (P-F40h), with $D/2=28$, proves
\[
                              C_1<0<C_3.
 \tag{P-F40k}
\]
This establishes mixed local signs even when inverse pairing
does not cancel a simple pole.
\end{proof}

\paragraph{Three fixed variables remove the odd-codimension pair cancellation.}
There is also a precise limit to the parity shortcut. Let
$n\ge11$ with $3\nmid n$, put $m=n-3$, and consider the smooth
cubic in $n$ variables
\[
             F=K_m(x_0,\ldots,x_{m-1})+y_1^3+y_2^3+y_3^3.
 \tag{P-F40l}
\]
Extend the preceding Klein automorphism on the $x$ variables by
the identity on the three $y$ variables, and denote it by
$\widetilde\beta$. It is in $SL_n$ and has order
$P_m=(2^m-(-1)^m)/3$. Its normal representation
$\widetilde S=W_n/(\mathfrak{sl}_n\cdot F)$ satisfies
\[
 \dim\widetilde S=D+1,\qquad
 \dim\widetilde S^{\widetilde\beta}=2,
 \qquad \det(\widetilde\beta\mid\widetilde S)=1.
 \tag{P-F40m}
\]
Here and below $D=\binom n3$ uses the total number of variables.
To prove the fixed-space assertion, first note that no two
Klein eigenvalues have product one when $m\ge5$ and $3\nmid m$.
Equal indices would require the odd integer $P_m>1$ to divide
a power of two. For distinct indices, after removing a power
of $-2$, such a pair would require
$P_m\mid1+(-2)^d$ with $1\le d\le m-1$.
For $d\le m-2$ its nonzero absolute value is smaller than
$P_m$, since $2^{m-2}+1<P_m$ for $m\ge5$.
For $d=m-1$ its absolute value lies strictly between
$P_m$ and $2P_m$; equality with $P_m$ would give
$2^{m-1}=2$ when $m$ is even and $2^{m-1}=-2$ when $m$
is odd. Neither is possible. Thus there are no such pairs.

No Klein eigenvalue itself is one. Consequently no mixed
cubic monomial involving both the $x$ and $y$ variables is
fixed: its weight is either one Klein weight or the sum of
two. The fixed cubic subspace has dimension
$m+\binom{3+2}{3}=m+10$.
The centralizer in $\mathfrak{sl}_n$ has dimension
$m+3^2-1=m+8$, because the Klein eigenvalues are distinct
and different from the eigenvalue one on the three fixed
variables. Their difference is two. Smoothness follows
from the independent gradient blocks, and the same regular
partials argument gives Lie-orbit injectivity. This proves
the dimension assertions in (P-F40m); its determinant
assertion follows from the same cubic/adjoint determinant
calculation as before.

Put
\[
 \widetilde C_a=
 \frac1{\det(1-\widetilde\beta^a\mid
                         \widetilde S^*_{\rm mov})}
                \quad((a,P_m)=1).
\]
The moving codimension is now $D-1$, so
$\widetilde C_{-a}=(-1)^{D-1}\widetilde C_a$.
In particular, if $n\equiv3\pmod4$, the inverse pair in the
cyclic normal-model Molien sum has normalized double-pole
coefficient
\[
                    \frac{2\widetilde C_a}{P_m}\ne0
                    \quad\hbox{at }t=1.
 \tag{P-F40n}
\]
Indeed each term is $(1-t)^{-2}$ times a holomorphic function
with value $\widetilde C_a$ at one. The two values are equal
and nonzero in this parity, so this particular pair has an
actual pole of order two. This does not exclude cancellation
against other pairs in the complete average.

In the standard embedding one again has
\[
                  \log_2|\widetilde C_1|
                             =\frac{n^4}{24}+O(n^3).
 \tag{P-F40o}
\]
To check this without a new estimate, use the bulk calculation
above on the first $m$ variables: its contribution is
$m^4/24+O(m^3)$. Cubic monomials involving one of the three
fixed variables number only $O(n^2)$, as do all adjoint
factors and all exceptional top-index factors. Each nontrivial
$P_m$th-root factor has logarithm $O(n)$, so these change the
answer by $O(n^3)$. Finally $m=n-3$ gives
$m^4/24=n^4/24+O(n^3)$.
Thus an odd value of $D$ in the full $n$-variable calculation does not remove every
large local contribution at one: a higher-dimensional fixed
normal space changes the relevant parity and Laurent order.

\paragraph{Scope of the local calculation and the later global comparison.}
The representation $S$ comes from a genuine stabilizer subgroup
of a smooth cubic. The full stabilizer may be larger than
$\Gamma$, and the cyclic model (P-F40g) is not claimed to be
the entire local quotient. Even a correctly identified local
quotient would not identify its ordinary tangent-coordinate
Hilbert grading with the global stretching grading of P-F31.
A global equivariant Riemann--Roch or Molien--Weyl comparison,
with all strata and their signs controlled, would be required
to express the actual $B_n$ through such local data.
P-F44 below supplies an exact global comparison for the
primitive-Klein collection by the Weyl--Brion expansion;
its normalization is $1/(nP)$ and it retains the entire
complementary signed contribution. The cyclic model itself,
and in particular the three-fixed-variable extension, is
not identified with the global Hilbert series. The exact
odd-codimension cancellation shows why summing absolute
local amplitudes cannot replace that comparison.

The classical family and its diagonal projective group are in
Zheng, \emph{On abelian automorphism groups of hypersurfaces},
Section~2.2 and Lemma~3.1, printed pp.~3 and~5,
\url{https://arxiv.org/pdf/2004.09008}.
The commutative-algebra premises are directly available at
\url{https://stacks.math.columbia.edu/tag/00NQ},
\url{https://stacks.math.columbia.edu/tag/02JN}, and
\url{https://stacks.math.columbia.edu/tag/062F}.
No historical priority is claimed for this family or for
Molien averaging; the phase and scope computations are
proved explicitly above.
